\documentclass[11pt]{article}

% pdfLaTeX-friendly (no fontspec)
\usepackage[margin=1in]{geometry}
\usepackage[T1]{fontenc}
\usepackage[utf8]{inputenc}
\usepackage{lmodern}

\usepackage{amsmath,amssymb}
\usepackage{xcolor}
\usepackage{hyperref}
\usepackage{listings}

% Dirac notation (keep \rangle correct)
\newcommand{\ket}[1]{\lvert #1 \rangle}
\newcommand{\bra}[1]{\langle #1 \rvert}

% Lean listings setup (handles Lean Unicode in code blocks)
\lstdefinelanguage{Lean}{
  morekeywords={import,open,scoped,namespace,section,end,variable,def,abbrev,lemma,theorem,by,fun,have,calc,refine,intro,exact,classical,simp,funext,by_cases,omit},
  sensitive=true,
  morecomment=[l]{--},
  morecomment=[s]{/-}{-/},
  morestring=[b]"
}

\lstset{
  language=Lean,
  basicstyle=\ttfamily\footnotesize,
  columns=fullflexible,
  breaklines=true,
  frame=single,
  rulecolor=\color{black!25},
  keywordstyle=\color{blue!70!black},
  commentstyle=\color{green!40!black},
  stringstyle=\color{orange!60!black},
  showstringspaces=false,
  keepspaces=true,
  literate=
    {ℕ}{{\ensuremath{\mathbb{N}}}}1
    {ℤ}{{\ensuremath{\mathbb{Z}}}}1
    {ℚ}{{\ensuremath{\mathbb{Q}}}}1
    {𝕜}{{\ensuremath{\mathbb{k}}}}1
    {→}{{\ensuremath{\to}}}1
    {↔}{{\ensuremath{\leftrightarrow}}}1
    {∑}{{\ensuremath{\sum}}}1
    {⊆}{{\ensuremath{\subseteq}}}1
    {∈}{{\ensuremath{\in}}}1
    {∀}{{\ensuremath{\forall}}}1
    {∃}{{\ensuremath{\exists}}}1
    {∧}{{\ensuremath{\wedge}}}1
    {≤}{{\ensuremath{\le}}}1
    {≥}{{\ensuremath{\ge}}}1
    {≠}{{\ensuremath{\neq}}}1
    {↑}{{\ensuremath{\uparrow}}}1
    {·}{{\ensuremath{\bullet}}}1
    {⦃}{{\{} }1
    {⦄}{{\}} }1
    {“}{{``}}1
    {”}{{''}}1
}

\title{Explanation of the Lean 4 (mathlib) Code for the Subset-Sum (SS) Condition}
\author{}
\date{}

\begin{document}
\maketitle

\section*{Big picture}

You're formalizing the ``subset-sum support constraint'' for transversal diagonal gates:

\begin{itemize}
\item A transversal diagonal gate assigns each computational basis string $s\in\{0,1\}^n$ a phase depending on a modular weighted sum
\[
\sum_j a_j s_j \bmod m.
\]
\item For a logical basis state
\[
\ket{k_L}=\sum_s c^{(k)}_s \ket{s}
\]
to be an eigenvector with eigenvalue $e^{2\pi i b_k/m}$, every string $s$ that appears with nonzero coefficient must satisfy
\[
\sum_j a_j s_j \equiv b_k \pmod m.
\]
\item You formalize that as: \textbf{support $\subseteq S_k$} (for finitary ``support sets'') or as: \textbf{if amplitude is nonzero then weight $=b_k$} (for functional amplitude vectors).
\end{itemize}

Then you prove a correctness lemma:

\begin{itemize}
\item If the SS condition holds, then the abstract diagonal action \texttt{Dact} acts as a \textbf{global phase} on the state.
\end{itemize}

This is the core ``SS $\Rightarrow$ eigenstate under diagonal action'' statement.

\section*{Imports and namespace}

\begin{lstlisting}
import Mathlib.Data.ZMod.Basic
import Mathlib.Data.Bool.Basic
import Mathlib.Data.Finset.Lattice.Basic
import Mathlib.Algebra.BigOperators.Fin
import Mathlib.Data.Rat.Defs
import Mathlib.Data.Fintype.Sets

open scoped BigOperators
namespace SS
\end{lstlisting}

\begin{itemize}
\item \texttt{ZMod.Basic}: gives \texttt{ZMod m} and basic algebra on it.
\item \texttt{Bool.Basic}: for \texttt{Bool} and simp lemmas about \texttt{if}.
\item \texttt{Finset.Lattice.Basic}: gives coercions \texttt{(↑supp : Set α)}, subset relations, etc.
\item \texttt{BigOperators.Fin}: provides \texttt{∑ i : Fin n, ...} notation.
\item \texttt{Rat.Defs}: provides the rational numbers \texttt{ℚ} with the additive/ordered structure needed for big sums.
\item \texttt{Fintype.Sets}: provides \texttt{Set.toFinset}, used to turn residue classes \texttt{Sk a (S j)} into \texttt{Finset}s.
\end{itemize}

All of your definitions are placed in \texttt{namespace SS} so names don't pollute globally.

\section*{BitString}

\begin{lstlisting}
abbrev BitString (n : ℕ) : Type := Fin n → Bool
\end{lstlisting}

This is your encoding of a binary string of length $n$.

\begin{itemize}
\item A string $s\in\{0,1\}^n$ is a function \texttt{Fin n → Bool}.
\item This is convenient in Lean because \texttt{Fin n} is the index type $\{0,\dots,n-1\}$.
\end{itemize}

\subsection*{Bitwise complement \texttt{complement}}

The BD$_{2m}$ symmetry constraint in the paper involves ``bitwise complement'' of strings.
In Lean this is implemented as Boolean negation at each coordinate.

\begin{lstlisting}
def complement {n : ℕ} (s : BitString n) : BitString n := fun i => !s i

@[simp] lemma complement_apply {n : ℕ} (s : BitString n) (i : Fin n) :
    complement s i = !s i := rfl

@[simp] lemma complement_complement {n : ℕ} (s : BitString n) :
    complement (complement s) = s := by
  funext i
  simp [complement]
\end{lstlisting}

So if $s\in\{0,1\}^n$, then $\overline{s}$ is defined by $(\overline{s})_i = 1 - s_i$.
The lemma \texttt{complement\_complement} formalizes that complement is an involution.

\section*{Section Basic}

\subsection*{Variables}

\begin{lstlisting}
section Basic
variable {m n : ℕ} [NeZero m]
\end{lstlisting}

Here you initially assume $m\neq 0$ so you can treat \texttt{ZMod m} as the usual modular ring.
You later discovered some lemmas don't need it, so you used \texttt{omit} for those.

\subsection*{\texttt{bit : Bool → ZMod m}}

\begin{lstlisting}
def bit (b : Bool) : ZMod m := if b then (1 : ZMod m) else 0
\end{lstlisting}

This is the map $\{0,1\}\to\mathbb{Z}_m$ that sends
\begin{itemize}
\item \texttt{false} $\mapsto 0$
\item \texttt{true} $\mapsto 1$
\end{itemize}
This is the bridge between your Boolean bitstrings and arithmetic in \texttt{ZMod m}.

\subsection*{simp lemmas \texttt{bit\_false}, \texttt{bit\_true}}

\begin{lstlisting}
omit [NeZero m] in
@[simp] lemma bit_false : bit (m := m) false = 0 := by
  simp [bit]

omit [NeZero m] in
@[simp] lemma bit_true : bit (m := m) true = 1 := by
  simp [bit]
\end{lstlisting}

They just tell \texttt{simp} how to reduce \texttt{bit false} and \texttt{bit true}.
You wrapped them in \texttt{omit [NeZero m]} because they don't actually use $m\neq 0$.

\subsection*{\texttt{weight}}

\begin{lstlisting}
def weight (a : Fin n → ZMod m) (s : BitString n) : ZMod m :=
  ∑ i : Fin n, a i * bit (m := m) (s i)
\end{lstlisting}

This is exactly the modular subset-sum:
\begin{itemize}
\item You provide an ``angle vector'' $a=(a_0,\dots,a_{n-1})$ as \texttt{Fin n → ZMod m}.
\item For a bitstring $s$, compute
\[
w_a(s)=\sum_i a_i\cdot s_i \pmod m,
\]
where $s_i\in\{0,1\}$ is implemented as \texttt{bit (s i)}.
\end{itemize}

This \texttt{weight} is the mathematical core of SS: it tells you what phase a diagonal gate assigns to $\ket{s}$.

\subsection*{\texttt{Sk} (support subset)}

\begin{lstlisting}
def Sk (a : Fin n → ZMod m) (bk : ZMod m) : Set (BitString n) :=
  { s | weight (m := m) a s = bk }
\end{lstlisting}

This is the SS subset
\[
S_k=\{\, s\in\{0,1\}^n : w_a(s)=b_k \}.
\]
In other words, \texttt{Sk a bk} is the set of computational basis strings that pick up the same phase label \texttt{bk}.

\subsection*{\texttt{mem\_Sk}}

\begin{lstlisting}
omit [NeZero m] in
@[simp] lemma mem_Sk {m n : ℕ} (a : Fin n → ZMod m) (bk : ZMod m) (s : BitString n) :
    s ∈ Sk (m := m) a bk ↔ weight (m := m) a s = bk :=
  Iff.rfl
\end{lstlisting}

Just the definitional unfolding of membership in that set. Again it doesn't need \texttt{NeZero}.

\subsection*{\texttt{SSConditionSupport}}

\begin{lstlisting}
def SSConditionSupport (a : Fin n → ZMod m) (bk : ZMod m) (supp : Finset (BitString n)) : Prop :=
  (↑supp : Set (BitString n)) ⊆ Sk (m := m) a bk
\end{lstlisting}

This is the SS condition stated as a support containment statement:

\begin{quote}
Every string in the (finite) support \texttt{supp} lies inside \texttt{Sk}.
\end{quote}

This matches the paper's statement $c^{(k)}_s\neq 0 \Rightarrow s\in S_k$,
but phrased for a \texttt{Finset} support rather than amplitudes.

\subsection*{\texttt{SSConditionSupport'}}

\begin{lstlisting}
def SSConditionSupport' (a : Fin n → ZMod m) (bk : ZMod m) (supp : Finset (BitString n)) : Prop :=
  ∀ ⦃s⦄, s ∈ supp → weight (m := m) a s = bk
\end{lstlisting}

Same condition, but in ``elementwise'' form:

\begin{quote}
If $s\in\texttt{supp}$, then its weight equals \texttt{bk}.
\end{quote}

\subsection*{\texttt{SSConditionSupport\_iff}}

\begin{lstlisting}
omit [NeZero m] in
theorem SSConditionSupport_iff
    {m n : ℕ} (a : Fin n → ZMod m) (bk : ZMod m) (supp : Finset (BitString n)) :
    SSConditionSupport (m := m) a bk supp ↔ SSConditionSupport' (m := m) a bk supp := by
  constructor
  · intro h s hs
    have : s ∈ (↑supp : Set (BitString n)) := by simpa using hs
    have : s ∈ Sk (m := m) a bk := h this
    simpa [Sk] using this
  · intro h s hs
    have hs' : s ∈ supp := by simpa using hs
    simpa [Sk] using h hs'

end Basic
\end{lstlisting}

This theorem proves those two formulations are logically equivalent:
\[
\texttt{SSConditionSupport a bk supp} \leftrightarrow  \texttt{SSConditionSupport' a bk supp}.
\]
So you can use whichever is more convenient when proving things.

\begin{itemize}
\item The $\Rightarrow$ direction: from set inclusion you get the weight equation using \texttt{Sk}'s definition.
\item The $\Leftarrow$ direction: from the pointwise property you build the set inclusion.
\end{itemize}

Again, you \texttt{omit [NeZero m]} because this proof doesn't need $m\neq 0$.

\section*{Section BDSymmetry (BD$_{2m}$ constraint)}

This is the new piece you added to match the BD$_{2m}$ (binary dihedral) symmetry constraint from the paper.

Physically, when the code has a transversal Pauli-$X$ of the form $X^{\otimes n}$ acting as a logical $X$, the logical basis states satisfy
\[
\ket{1_L} = X^{\otimes n} \ket{0_L}.
\]
Applying $X$ flips every physical bit, so at the level of basis strings and supports you expect a ``complement symmetry'':
\[
S_1 = \{\, \overline{s} : s \in S_0 \}.
\]

In the paper, BD$_{2m}$ symmetry also enforces a constraint on the angle vector $a$:
\[
\sum_i a_i \equiv -1 \pmod m,
\]
which ensures $0^n \in S_0$ and $1^n \in S_1$.

\subsection*{Angle constraint \texttt{BDAngleCondition}}

\begin{lstlisting}
section BDSymmetry

variable {m n : ℕ}

/-- BD_{2m} angle constraint: the total angle sum is `-1` in `ZMod m`. -/
def BDAngleCondition (a : Fin n → ZMod m) : Prop :=
  ∑ i : Fin n, a i = -1
\end{lstlisting}

This is a direct translation of $\sum_i a_i \equiv -1 \pmod m$ into Lean.
(It does not require the assumption \texttt{[NeZero m]}.)

\subsection*{A helper simp lemma: \texttt{bit\_not}}

To relate weights of $s$ and $\overline{s}$, you first prove that Boolean negation corresponds to $1-\texttt{bit}$:

\begin{lstlisting}
@[simp] lemma bit_not {m : ℕ} (b : Bool) :
    bit (m := m) (!b) = 1 - bit (m := m) b := by
  cases b <;> simp [bit]
\end{lstlisting}

Mathematically, this is just:
\[
\mathrm{bit}(1-b) = 1 - \mathrm{bit}(b).
\]

\subsection*{Key identity: weight of a complemented string}

\begin{lstlisting}
theorem weight_complement (a : Fin n → ZMod m) (s : BitString n) :
    weight (m := m) a (complement s) =
      (∑ i : Fin n, a i) - weight (m := m) a s := by
  classical
  simp [weight, complement, bit_not, mul_sub, Finset.sum_sub_distrib]
\end{lstlisting}

This is the formal version of the computation
\[
\sum_i a_i (1-s_i) = \left(\sum_i a_i\right) - \sum_i a_i s_i.
\]

\subsection*{Main lemma: \texttt{complement\_mem\_S1\_of\_mem\_S0}}

With the angle constraint and the SS definition of $S_k$, you prove:

\begin{lstlisting}
theorem complement_mem_S1_of_mem_S0
    (a : Fin n → ZMod m) (hBD : BDAngleCondition (m := m) a)
    (s : BitString n) (hs : s ∈ Sk (m := m) a 0) :
    complement s ∈ Sk (m := m) a (-1) := by
  simp [Sk] at hs ⊢
  have hsum : (∑ i : Fin n, a i) = (-1 : ZMod m) := by
    simpa [BDAngleCondition] using hBD
  calc
    weight (m := m) a (complement s)
        = (∑ i : Fin n, a i) - weight (m := m) a s := weight_complement (m := m) a s
    _ = (-1 : ZMod m) - 0 := by simp [hsum, hs]
    _ = (-1 : ZMod m) := by simp
\end{lstlisting}

This is exactly the paper's reasoning:
if $w_a(s)=0$ and $\sum_i a_i = -1$, then
\[
w_a(\overline{s}) = \left(\sum_i a_i\right) - w_a(s) = -1 - 0 = -1.
\]

For the set equality later, you also prove the reverse direction
\texttt{complement\_mem\_S0\_of\_mem\_S1}, using that complement is an involution.

\subsection*{Supports are complement-related: \texttt{Sk\_complement\_eq}}

\begin{lstlisting}
theorem Sk_complement_eq (a : Fin n → ZMod m) (hBD : BDAngleCondition (m := m) a) :
    Sk (m := m) a (-1) = complement '' Sk (m := m) a 0 := by
  ext s
  constructor
  · intro hs
    refine ⟨complement s, complement_mem_S0_of_mem_S1 (m := m) a hBD s hs, ?_⟩
    simp
  · rintro ⟨t, ht, rfl⟩
    exact complement_mem_S1_of_mem_S0 (m := m) a hBD t ht
\end{lstlisting}

This says: under BD$_{2m}$ symmetry, the phase-$(-1)$ SS set is exactly the image of the phase-$0$ SS set under bitwise complement.

\section*{Section ResidueShiftScreen (Classical Distance Condition)}

This section formalizes the \textbf{Residue-Shift Screen} from paper 2510.20728 (Equation 3).
The screen condition ensures that no two codewords from different residue classes are Hamming-1 neighbors,
which guarantees the classical distance $d(C) \geq 2$.
This is crucial because it makes X-type and Y-type Knill-Laflamme conditions vanish automatically.

\subsection*{Bit flipping: \texttt{flipBit}}

First, we define the operation of flipping a single bit at position $i$:

\begin{lstlisting}
def flipBit (s : BitString n) (i : Fin n) : BitString n :=
  fun j => if j = i then !s j else s j
\end{lstlisting}

This creates a Hamming-1 neighbor of $s$ by toggling the $i$-th bit.
We also prove basic properties:

\begin{lstlisting}
@[simp] lemma flipBit_apply_eq (s : BitString n) (i : Fin n) :
    flipBit s i i = !s i := by simp [flipBit]

@[simp] lemma flipBit_apply_ne (s : BitString n) (i j : Fin n) (hij : j ≠ i) :
    flipBit s i j = s j := by simp [flipBit, hij]

@[simp] lemma flipBit_flipBit (s : BitString n) (i : Fin n) :
    flipBit (flipBit s i) i = s
\end{lstlisting}

The last lemma shows that flipping the same bit twice returns the original string.

\subsection*{Key lemma: \texttt{weight\_flipBit}}

The central technical result relates the weight of a flipped string to the original:

\begin{lstlisting}
theorem weight_flipBit (a : Fin n → ZMod m) (s : BitString n) (i : Fin n) :
    weight a (flipBit s i) = weight a s + (if s i then -a i else a i)
\end{lstlisting}

Mathematically, this says:
\[
w_a(\text{flip}_i(s)) = w_a(s) + \begin{cases} -a_i & \text{if } s_i = 1 \\ +a_i & \text{if } s_i = 0 \end{cases}
\]

The proof isolates the $i$-th term in the sum using \texttt{Finset.sum\_eq\_single\_of\_mem}:
\begin{itemize}
\item For $j \neq i$: the $j$-th term is unchanged since $(\text{flip}_i(s))_j = s_j$.
\item For $j = i$: the term changes from $a_i \cdot s_i$ to $a_i \cdot (1 - s_i)$, giving a difference of $\pm a_i$.
\end{itemize}

We also provide a variant expressing this as a difference:

\begin{lstlisting}
theorem weight_flipBit_sub (a : Fin n → ZMod m) (s : BitString n) (i : Fin n) :
    weight a (flipBit s i) - weight a s = if s i then -a i else a i
\end{lstlisting}

\subsection*{The screen condition: \texttt{ResidueShiftScreen}}

\begin{lstlisting}
def ResidueShiftScreen (a : Fin n → ZMod m) (S : Fin K → ZMod m) : Prop :=
  ∀ i : Fin n, ∀ j k : Fin K, j ≠ k →
    S j - S k ≠ a i ∧ S j - S k ≠ -a i
\end{lstlisting}

This is the formal statement of Equation (3) from the paper:
\[
S_j - S_k \not\equiv \pm a_i \pmod{m} \quad \text{for all } i \text{ and } j \neq k.
\]

Here $S : \texttt{Fin K} \to \texttt{ZMod m}$ assigns phase labels to the $K$ logical basis states.
The condition forbids the difference between any two distinct phase labels from equaling $\pm a_i$ for any weight component $a_i$.

\subsection*{Main theorem: \texttt{no\_hamming1\_neighbor\_of\_screen}}

\begin{lstlisting}
theorem no_hamming1_neighbor_of_screen
    (a : Fin n → ZMod m) (S : Fin K → ZMod m)
    (hScreen : ResidueShiftScreen a S)
    {j k : Fin K} (hjk : j ≠ k)
    {s : BitString n} (hs : s ∈ Sk a (S j))
    {i : Fin n} :
    flipBit s i ∉ Sk a (S k)
\end{lstlisting}

This is the key result: \emph{if the screen condition holds and $s \in S_j$, then flipping any bit of $s$ cannot land in $S_k$ for $k \neq j$}.

The proof proceeds by contradiction:
\begin{enumerate}
\item Assume $\text{flip}_i(s) \in S_k$.
\item From \texttt{weight\_flipBit}: $w_a(\text{flip}_i(s)) = w_a(s) + (\pm a_i) = S_j + (\pm a_i)$.
\item But $\text{flip}_i(s) \in S_k$ means $w_a(\text{flip}_i(s)) = S_k$.
\item Therefore $S_k - S_j = \pm a_i$, contradicting the screen condition.
\end{enumerate}

\subsection*{Corollaries}

We provide two additional formulations:

\begin{lstlisting}
theorem disjoint_hamming1_of_screen
    (a : Fin n → ZMod m) (S : Fin K → ZMod m)
    (hScreen : ResidueShiftScreen a S) :
    ∀ j k : Fin K, j ≠ k → ∀ s ∈ Sk a (S j), ∀ i : Fin n, flipBit s i ∉ Sk a (S k)
\end{lstlisting}

This restates the main theorem in a more explicit ``for all'' form:
the union of all supports has no Hamming-1 neighbors between different classes.

\begin{lstlisting}
theorem hamming_distance_ge_2_of_screen
    (a : Fin n → ZMod m) (S : Fin K → ZMod m)
    (hScreen : ResidueShiftScreen a S)
    {j k : Fin K} (hjk : j ≠ k)
    {s t : BitString n} (hs : s ∈ Sk a (S j)) (ht : t ∈ Sk a (S k)) :
    ¬ ∃ i : Fin n, t = flipBit s i
\end{lstlisting}

This is the distance formulation: any two codewords from different residue classes have Hamming distance $\geq 2$.
In other words, no codeword in $S_k$ can be obtained by flipping a single bit of any codeword in $S_j$ (for $j \neq k$).

\subsection*{Physical significance}

The Residue-Shift Screen is important because:
\begin{itemize}
\item It guarantees the classical code (union of all $S_k$) has minimum distance $d \geq 2$.
\item This makes single-qubit X and Y errors always detectable.
\item In the Knill-Laflamme framework, X-type and Y-type conditions become trivially satisfied when $d \geq 2$, since such errors cannot map between valid codewords.
\item Only Z-type Knill--Laflamme conditions need to be checked separately; these are formalized in Section~ZTypeKL below as linear constraints on probability weights.
\end{itemize}

\section*{Section Disjointness}

\subsection*{Setup}

\begin{lstlisting}
section Disjointness

variable {m n K : ℕ}      -- removed [NeZero m]
variable (a : Fin n → ZMod m)
variable (b : Fin K → ZMod m)
variable (supp : Fin K → Finset (BitString n))
\end{lstlisting}

Here you introduce:
\begin{itemize}
\item \texttt{b : Fin K → ZMod m} = the list of intended logical phase labels $(b_0,\dots,b_{K-1})$.
\item \texttt{supp k} = the finite support set for the $k$-th logical basis state.
\end{itemize}

\subsection*{Theorem \texttt{supports\_disjoint\_of\_injective}}

\begin{lstlisting}
theorem supports_disjoint_of_injective
    (hb : Function.Injective b)
    (hSS : ∀ k, SSConditionSupport' (m := m) a (b k) (supp k)) :
    ∀ {k l : Fin K}, k ≠ l → Disjoint (supp k) (supp l) := by
  intro k l hkl
  classical
  refine Finset.disjoint_left.2 ?_
  intro s hsk hsl
  have hk : weight (m := m) a s = b k := hSS k hsk
  have hl : weight (m := m) a s = b l := hSS l hsl
  have hbl : b k = b l := by
    calc
      b k = weight (m := m) a s := hk.symm
      _   = b l               := hl
  exact hkl (hb hbl)

end Disjointness
\end{lstlisting}

Meaning:
\begin{itemize}
\item If the phase labels \texttt{b k} are all distinct (injective), and
\item each support \texttt{supp k} satisfies SS with its label \texttt{b k},
\item then supports for different logical indices cannot overlap.
\end{itemize}

This is the orthogonality-by-disjoint-support fact: if $b_k\neq b_\ell$, no basis string can appear in both
$\ket{k_L}$ and $\ket{\ell_L}$ under the SS rule.

How the proof works (in plain math):
\begin{itemize}
\item Suppose a string $s$ lies in both \texttt{supp k} and \texttt{supp l}.
\item Then SS says \texttt{weight a s = b k} and also \texttt{weight a s = b l}.
\item Hence \texttt{b k = b l}, contradicting injectivity unless $k=l$.
\item Therefore they are disjoint.
\end{itemize}

This matches the statement in your manuscript that distinct $b_k$ make $S_k$ disjoint and thus (at least) support-level orthogonality.

\section*{Section DiagonalAction}

This section switches from ``supports as finsets'' to ``full amplitude vectors''.

\subsection*{Variables}

\begin{lstlisting}
section DiagonalAction

variable {m n : ℕ} -- no [NeZero m]
variable {𝕜 : Type*} [MonoidWithZero 𝕜]
\end{lstlisting}

\begin{itemize}
\item \texttt{𝕜} is the coefficient type for amplitudes. You only assume \texttt{MonoidWithZero} because you need multiplication (\texttt{*}) and a zero element.
\item You deliberately don't assume \texttt{𝕜 = ℂ} yet: this keeps it abstract and reusable.
\end{itemize}

\subsection*{\texttt{State}}

\begin{lstlisting}
abbrev State (n : ℕ) (k : Type*) : Type _ := BitString n → k
\end{lstlisting}

A ``state'' is an amplitude assignment $\psi:\{0,1\}^n\to\mathbb{k}$.
So in physics notation: $\psi(s)$ is the coefficient of $\ket{s}$.
(You fixed the universe issue by writing \texttt{Type \_}.)

\subsection*{\texttt{SSConditionState}}

\begin{lstlisting}
def SSConditionState (a : Fin n → ZMod m) (bk : ZMod m) (ψ : State n 𝕜) : Prop :=
  ∀ s, ψ s ≠ 0 → weight (m := m) a s = bk
\end{lstlisting}

This is the most faithful Lean version of the manuscript SS condition:
if amplitude at $s$ is nonzero, then $s$ must have the right subset-sum.
It is literally the ``support lies in $S_k$'' condition.
\subsection*{\texttt{Dact}}

\begin{lstlisting}
def Dact (phase : ZMod m → 𝕜) (a : Fin n → ZMod m) (ψ : State n 𝕜) : State n 𝕜 :=
  fun s => phase (weight (m := m) a s) * ψ s
\end{lstlisting}

This is an abstract diagonal operator acting on amplitude functions.

\begin{itemize}
\item In the computational basis, a diagonal gate acts by multiplying each basis vector $\ket{s}$ by some scalar depending on $s$.
\item Here, the scalar is \texttt{phase (weight a s)}.
\end{itemize}

So this corresponds exactly to the transversal diagonal gate's induced action
\[
(D\psi)(s)=\omega^{w_a(s)} \psi(s),
\]
but abstracted as \texttt{phase : ZMod m → 𝕜} instead of hardcoding $\omega^x$.

\subsection*{Theorem \texttt{Dact\_eq\_global\_of\_SS}}

\begin{lstlisting}
theorem Dact_eq_global_of_SS
    (phase : ZMod m → 𝕜) (a : Fin n → ZMod m) (bk : ZMod m) (ψ : State n 𝕜)
    (hSS : SSConditionState (m := m) a bk ψ) :
    Dact (m := m) phase a ψ = fun s => phase bk * ψ s := by
  classical
  funext s
  by_cases hs : ψ s = 0
  · simp [Dact, hs]
  · have hw : weight (m := m) a s = bk := hSS s hs
    simp [Dact, hw]

end DiagonalAction
\end{lstlisting}

This is the ``SS correctness'' lemma:
If $\psi$ only has support on strings with weight $=b_k$, then applying the diagonal action multiplies every nonzero component by the same scalar \texttt{phase bk}.
So the whole state is an eigenvector with eigenvalue \texttt{phase bk}.

How the proof works:
\begin{itemize}
\item Prove equality of functions, so \texttt{funext s}.
\item Split into two cases:
\begin{enumerate}
\item $\psi(s)=0$: both sides are $0$ after multiplication, so \texttt{simp} closes it.
\item $\psi(s)\neq 0$: SS tells you \texttt{weight a s = bk}, so the phase is \texttt{phase bk}, and \texttt{simp} closes it.
\end{enumerate}
\end{itemize}

This is the formal Lean statement that SS implies the state is an eigenstate of the diagonal gate.

\section*{Section ZTypeKL (Z-type Knill--Laflamme conditions)}

This section formalizes the \textbf{Z-type Knill--Laflamme (KL)} constraints from paper 2510.20728 (Eq.~6).
For each physical qubit $i$, the expectation value of $Z_i$ must be the same across all logical states.

If a logical state $\ket{j_L}$ has a probability distribution $p_{j,x}$ supported on some finite set \texttt{supp j},
the condition is:
\[
\langle j_L|Z_i|j_L\rangle = \sum_{x \in \text{supp}_j} p_{j,x}\,(1-2x_i) = t_i,
\]
where $t_i$ is independent of $j$.

\subsection*{Sign for Z measurement: \texttt{zSign}}

\begin{lstlisting}
def zSign (s : BitString n) (i : Fin n) : ℤ := if s i then -1 else 1
\end{lstlisting}

This encodes $(-1)^{x_i}$ (equivalently $1-2x_i$).

\subsection*{Z expectation over a finite support: \texttt{zExpectation}}

\begin{lstlisting}
def zExpectation (supp : Finset (BitString n)) (p : BitString n → ℚ) (i : Fin n) : ℚ :=
  ∑ x ∈ supp, p x * zSign x i
\end{lstlisting}

\subsection*{Z-type KL constraints across $K$ logical states}

\begin{lstlisting}
def ZTypeKL (K : ℕ) (supp : Fin K → Finset (BitString n))
    (prob : Fin K → BitString n → ℚ) : Prop :=
  ∀ i : Fin n, ∀ j k : Fin K,
    zExpectation (supp j) (prob j) i = zExpectation (supp k) (prob k) i
\end{lstlisting}

\subsection*{Alternative formulation with explicit targets}

\begin{lstlisting}
def ZTypeKL' (K : ℕ) (supp : Fin K → Finset (BitString n))
    (prob : Fin K → BitString n → ℚ) (target : Fin n → ℚ) : Prop :=
  ∀ j : Fin K, ∀ i : Fin n, zExpectation (supp j) (prob j) i = target i
\end{lstlisting}

\subsection*{Equivalence theorem}

\begin{lstlisting}
theorem ZTypeKL_iff_exists_target (K : ℕ) [NeZero K]
    (supp : Fin K → Finset (BitString n)) (prob : Fin K → BitString n → ℚ) :
    ZTypeKL (n := n) K supp prob ↔ ∃ target : Fin n → ℚ, ZTypeKL' (n := n) K supp prob target := by
  constructor
  · intro h
    have hK : 0 < K := Nat.pos_of_ne_zero (NeZero.ne K)
    let j0 : Fin K := ⟨0, hK⟩
    refine ⟨fun i => zExpectation (supp j0) (prob j0) i, ?_⟩
    intro j i
    simpa using h i j j0
  · rintro ⟨target, htarget⟩
    intro i j k
    calc
      zExpectation (supp j) (prob j) i = target i := htarget j i
      _ = zExpectation (supp k) (prob k) i := (htarget k i).symm
\end{lstlisting}

The forward direction simply picks a reference logical index $j_0$ (using \texttt{[NeZero K]}) and defines the target values as the expectations for that state.

\subsection*{Normalization of probability distributions}

\begin{lstlisting}
def IsNormalizedProb (supp : Finset (BitString n)) (p : BitString n → ℚ) : Prop :=
  (∀ x, x ∉ supp → p x = 0) ∧ (∀ x ∈ supp, 0 ≤ p x) ∧ (∑ x ∈ supp, p x = 1)
\end{lstlisting}

\subsection*{Connecting to residue classes: \texttt{SSLPSetup}}

In the full SS-LP setup, supports are the residue classes \texttt{Sk a (S j)} turned into finsets using \texttt{Set.toFinset}.

\begin{lstlisting}
def SSLPSetup (n K : ℕ) (m : ℕ) [NeZero m]
    (a : Fin n → ZMod m) (S : Fin K → ZMod m)
    (prob : Fin K → BitString n → ℚ) : Prop :=
  by
    classical
    let supp : Fin K → Finset (BitString n) := fun j => (Sk (m := m) a (S j)).toFinset
    exact (∀ j, IsNormalizedProb (n := n) (supp j) (prob j)) ∧ ZTypeKL (n := n) K supp prob
\end{lstlisting}

\section*{What this file verifies (and what it doesn't yet)}

\subsection*{Verified / formalized}

\begin{itemize}
\item The definition of subset-sum weights on bitstrings.
\item The definition of the SS subsets $S_k$.
\item The SS condition as ``support subset'' and as ``nonzero amplitude $\Rightarrow$ correct congruence''.
\item Bitwise complement on bitstrings (\texttt{complement}) and the involution property \texttt{complement\_complement}.
\item The BD$_{2m}$ symmetry constraint on angles (\texttt{BDAngleCondition}) and its consequences, including \texttt{complement\_mem\_S1\_of\_mem\_S0} and \texttt{Sk\_complement\_eq}.
\item \textbf{Residue-Shift Screen} (Eq.~3 of paper 2510.20728):
  \begin{itemize}
  \item Bit-flip operation \texttt{flipBit} and its properties.
  \item Key lemma \texttt{weight\_flipBit}: weight changes by $\pm a_i$ under bit flip.
  \item Screen condition \texttt{ResidueShiftScreen}: $S_j - S_k \not\equiv \pm a_i \pmod{m}$.
  \item Main theorem \texttt{no\_hamming1\_neighbor\_of\_screen}: screen $\Rightarrow$ no Hamming-1 neighbors between classes.
  \item Corollary \texttt{hamming\_distance\_ge\_2\_of\_screen}: classical distance $d(C) \geq 2$.
  \end{itemize}
\item SS $\Rightarrow$ disjoint supports (under distinct phase labels).
\item SS $\Rightarrow$ eigenstate under the abstract diagonal action (global phase).
\end{itemize}

\subsection*{Not yet formalized}

\begin{itemize}
\item Z-type Knill-Laflamme conditions (linear constraints on amplitudes: $\sum_x (1-2x_i) p_{j,x} = t_i$).
\item Concrete code verification (e.g., verifying specific codes from the catalogue).
\item Control phases for multi-control gates.
\end{itemize}

\end{document}
